\textbf{Lösung zu Klausurvorbereitung 5}

\begin{loesung}
% KORRIGIERT: Original: "x + y = 2m für x \neq y, andernfalls spielt x gegen 2m" (passt nicht zur Tabelle).
In Runde $r$ spielt $x$ gegen $y$, wenn $x + y \equiv r \pmod{2m-1}$ und $x \neq y$ ist ($x, y \in \{1, \dots, 2m-1\}$); gilt $2x \equiv r \pmod{2m-1}$, so spielt $x$ gegen $2m$.

\medskip
\begin{tabular}{ccccccc}
\toprule
Runde 1 & Runde 2 & Runde 3 & Runde 4 & Runde 5 & Runde 6 & Runde 7\\
\midrule
1\ 7 & 1\ 8 & 1\ 2 & 1\ 3 & 1\ 4 & 1\ 5 & 1\ 6\\
2\ 6 & 2\ 7 & 3\ 7 & 2\ 8 & 2\ 3 & 2\ 4 & 2\ 5\\
3\ 5 & 3\ 6 & 4\ 6 & 4\ 7 & 5\ 7 & 3\ 8 & 3\ 4\\
4\ 8 & 4\ 5 & 5\ 8 & 5\ 6 & 6\ 8 & 6\ 7 & 7\ 8\\
\bottomrule
\end{tabular}
\end{loesung}
